% ============================================================================
% Section 1.1 — Unit Rates with Fractions
% CCSS 7.RP.A.1
% ============================================================================
\section{Unit Rates with Fractions}

\begin{stepGoal}
\begin{itemize}
    \item Compute unit rates when quantities are fractions
    \item Compare rates by finding each unit rate
\end{itemize}
\end{stepGoal}

\begin{stepVocabBox}
    \vocabItem{Unit Rate}{How much of one quantity for \textbf{one unit} of the other.}
    \vocabItem{Reciprocal}{The ``flipped'' fraction --- the reciprocal of $\frac{a}{b}$ is $\frac{b}{a}$.}
\end{stepVocabBox}

\begin{stepsCard}[How to Find a Unit Rate with Fractions]
    \stepItem{Write the rate as a \textbf{fraction}: $\dfrac{\text{first quantity}}{\text{second quantity}}$.}
    \stepItem{\textbf{Divide} by multiplying by the reciprocal (Keep--Change--Flip).}
    \stepItem{\textbf{Simplify} the result and label with units (``per one'').}
\end{stepsCard}

% --- Worked Example 1 ---
\begin{stepExample}{A recipe uses $\frac{3}{4}$ cup of sugar for $\frac{1}{2}$ batch. How much sugar per batch?}
    \stepShow{1} Write the rate: $\dfrac{\frac{3}{4} \text{ cup}}{\frac{1}{2} \text{ batch}}$

    \stepShow{2} Keep--Change--Flip: $\dfrac{3}{4} \times \dfrac{2}{1} = \dfrac{6}{4}$

    \stepShow{3} Simplify: $\dfrac{6}{4} = \dfrac{3}{2} = 1\dfrac{1}{2}$
    \stepResult{$1\dfrac{1}{2}$ cups per batch}
\end{stepExample}

% --- Worked Example 2 ---
\begin{stepExample}{Who walks faster? Ana: $\frac{2}{3}$ mi in $\frac{1}{4}$ hr. Ben: $\frac{3}{4}$ mi in $\frac{1}{3}$ hr.}
    \stepShow{1} Ana's rate: $\dfrac{\frac{2}{3}}{\frac{1}{4}}$ \qquad Ben's rate: $\dfrac{\frac{3}{4}}{\frac{1}{3}}$

    \stepShow{2} Ana: $\dfrac{2}{3} \times \dfrac{4}{1} = \dfrac{8}{3}$ \qquad Ben: $\dfrac{3}{4} \times \dfrac{3}{1} = \dfrac{9}{4}$

    \stepShow{3} Ana: $\dfrac{8}{3} = 2\dfrac{2}{3}$ mph \qquad Ben: $\dfrac{9}{4} = 2\dfrac{1}{4}$ mph
    \stepResult{Ana is faster ($2\frac{2}{3}$ mph $> 2\frac{1}{4}$ mph).}
\end{stepExample}

\mascotStepTip{``Per'' means ``for every one.'' Always divide to find a unit rate!}

% ============================================================================
% PRACTICE
% ============================================================================
\newpage
\begin{stepPractice}{Unit Rates with Fractions Practice}
    \resetProblems

    \stepPracticeHeader[step1Color]{Write the Rate}
    \begin{multicols}{2}
    \prob A painter covers $\frac{2}{3}$ wall in $\frac{1}{3}$ hr. Write as a fraction. \answerBlank[2.5cm]
    \answer{$\frac{2/3}{1/3}$}
    \prob A faucet leaks $\frac{1}{2}$ cup in $\frac{3}{4}$ min. Write as a fraction. \answerBlank[2.5cm]
    \answer{$\frac{1/2}{3/4}$}
    \end{multicols}

    \stepPracticeHeader[step2Color]{Divide and Simplify}
    \prob A painter covers $\frac{2}{3}$ wall in $\frac{1}{3}$ hr. How much wall per hour? \answerBlank[2.5cm]
    \answer{$2$ walls per hour}

    \prob A faucet leaks $\frac{1}{2}$ cup in $\frac{3}{4}$ min. How many cups per minute? \answerBlank[2.5cm]
    \answer{$\frac{2}{3}$ cup per minute}

    \stepPracticeHeader[step3Color]{Compare Unit Rates}
    \prob A snail crawls $\frac{5}{8}$ m in $\frac{3}{4}$ hr. A beetle crawls $\frac{7}{10}$ m in $\frac{4}{5}$ hr. Which is faster? \answerBlank[2cm]
    \answer{Beetle ($\frac{7}{8}$ m/hr $> \frac{5}{6}$ m/hr)}

    \wordProblem{A car uses $\frac{5}{6}$ gallon to travel $\frac{2}{3}$ mile. How many gallons per mile?}{gallons per mile}
    \answer{$1\frac{1}{4}$ gallons per mile}
\end{stepPractice}
